% !TeX program = pdflatex
% Compile com \documentclass[handout]{beamer} para a versão com as demonstrações completas.
\documentclass[10pt]{beamer}

\usetheme{default}
\usecolortheme{seagull}
\setbeamertemplate{navigation symbols}{}
\setbeamertemplate{theorems}[numbered]

\usepackage[T1]{fontenc}
\usepackage[utf8]{inputenc}
\usepackage[brazil]{babel}
\usepackage{amsmath,amssymb,amsthm}
\usepackage{stmaryrd}
\usepackage{tikz}

\setbeamertemplate{qed symbol}{$\blacksquare$}
\renewcommand{\qedsymbol}{$\blacksquare$}

% Rodapé com a seção corrente e o número do slide.
\setbeamertemplate{footline}{%
	\hfill{\scriptsize\usebeamercolor[fg]{structure}\insertsection\qquad\insertframenumber\,/\,\inserttotalframenumber}\hspace*{2ex}\vspace{1.2ex}%
}

% Pausas. \passo revela o bloco seguinte no próximo passo do frame.
% Ele usa o mesmo contador de <+->, então pode ser misturado com
% \uncover<+->{...} e \item<+-> sem numeração manual dos slides.
% \onslide<.-> mostra o que vem a seguir junto com o último passo.
\newcommand{\passo}{\onslide<+->}

\newcommand{\diag}[1]{\Delta_{#1}}
\newcommand{\inv}[1]{{#1}^{-1}}
\newcommand{\Fix}{\operatorname{Fix}}
\newcommand{\Img}{\operatorname{Img}}
\newcommand{\Dom}{\operatorname{Dom}}
\newcommand{\Cod}{\operatorname{Cod}}
% Composição em ordem de diagrama. Troque por \circ se preferir o símbolo usual.
\newcommand{\comp}{\mathbin{\fatsemi}}

\newcommand{\ida}{\textbf{($\Rightarrow$)}\ }
\newcommand{\volta}{\textbf{($\Leftarrow$)}\ }
\newcommand{\ordemnota}{%
	\smallskip
	{\footnotesize Lembrete: em $R \comp S$ aplica-se $R$ primeiro.}%
}

\newtheorem{prop}{Proposição}
\newtheorem{cor}{Corolário}

\newcommand{\relunit}{0.44cm}
\definecolor{relhlcolor}{RGB}{176,48,48}

\tikzset{
	relpic/.style   = {baseline=(current bounding box.center),
		x=\relunit, y=-\relunit},
	relcell/.style  = {draw=black!35, line width=0.3pt},
	reldot/.style   = {circle, fill=black!80, inner sep=0pt, minimum size=4.2pt},
	rellab/.style   = {font=\scriptsize, inner sep=1pt},
	relhl/.style    = {draw=relhlcolor, line width=0.9pt}
}

\newcommand{\reldraw}[2]{%
	\foreach \row [count=\i] in {#1} {%
		\xdef\relnrows{\i}%
		\foreach \e [count=\j] in \row {%
			\xdef\relncols{\j}%
			\draw[relcell] (\j-1,\i-1) rectangle (\j,\i);
			\ifnum\e>0 \node[reldot] at (\j-0.5,\i-0.5) {}; \fi
		}%
	}%
	#2%
}

\newcommand{\relmat}[4][]{%
	\begin{tikzpicture}[relpic]
		\reldraw{#4}{#1}
		\foreach \l [count=\i] in {#2} {\node[rellab,anchor=east]  at (-0.18,\i-0.5) {\l};}
		\foreach \l [count=\j] in {#3} {\node[rellab,anchor=north] at (\j-0.5,\relnrows+0.18) {\l};}
	\end{tikzpicture}%
}

\newcommand{\relgrid}[2][]{%
	\begin{tikzpicture}[relpic]
		\reldraw{#2}{#1}
	\end{tikzpicture}%
}

\newcommand{\hlrow}[1]{\draw[relhl] (0,#1-1) rectangle (\relncols,#1);}
\newcommand{\hlcol}[1]{\draw[relhl] (#1-1,0) rectangle (#1,\relnrows);}
\newcommand{\hlcell}[2]{\draw[relhl] (#2-1,#1-1) rectangle (#2,#1);}

\newcommand{\reldiag}{\draw[black!40,dashed,line width=0.3pt] (0,\relnrows) -- (\relncols,0);}

\newcommand{\relaxes}{%
	\node[rellab,anchor=south east] at (-0.10,-0.12) {$\Cod(R)$};
	\node[rellab,anchor=north west] at (\relncols+0.10,\relnrows+0.12) {$\Dom(R)$};
}

% Cadeia das quatro caixas encaixadas.
\newcommand{\cadeia}{%
	\begin{tikzpicture}[x=1cm,y=1cm,every node/.style={font=\scriptsize}]
		\draw[rounded corners] (0,0) rectangle (5.6,3.2);
		\node[anchor=north east] at (5.45,3.14) {$\Dom(R)$};
		\draw[rounded corners] (0.3,0.25) rectangle (5.3,2.6);
		\node[anchor=north east] at (5.15,2.54) {$\Img(R) \cap \Dom(R)$};
		\draw[rounded corners] (0.6,0.5) rectangle (5.0,2.0);
		\node[anchor=north east] at (4.85,1.94) {$\Fix(R \comp R)$};
		\draw[rounded corners] (0.9,0.75) rectangle (4.7,1.4);
		\node[anchor=north east] at (4.55,1.34) {$\Fix(R)$};
	\end{tikzpicture}%
}

\title{Entrega 8}
\date{}

\begin{document}
	
	\begin{frame}
		\titlepage
	\end{frame}
	
	\begin{frame}{Definições}
		\passo
		Fixe conjuntos $x$ e $y$ e uma relação $R \subseteq x \times y$. Escrevemos
		\[
		\Dom(R) := x, \qquad \Cod(R) := y,
		\]
		para o domínio e o contradomínio de $R$. Eles fazem parte da estrutura, pois o conjunto de pares $R$
		sozinho não os determina.
		\medskip
		
		\passo
		$R$ é chamada de \emph{função} quando, para todo $a \in \Dom(R)$, existe um único
		$b \in \Cod(R)$ com $aRb$.
		
		\medskip
		\passo
		\begin{align*}
			\diag{z} &:= \{(a,a) : a \in z\}, \\
			\inv{R}  &:= \{(b,a) : aRb\}, \\
			R \comp S &:= \{(a,c) : \exists b, \big(aRb \wedge bSc\big)\}, \\
			\Img(R) &:= \{b : \exists a, aRb\}, \\
			\Fix(R) &:= \{a : aRa\}.
		\end{align*}
		
		\medskip
		\passo
		O símbolo $\fatsemi$ marca a \textbf{ordem de diagrama}: em $R \comp S$ aplica-se $R$ primeiro.
		Em particular $\diag{\Dom(R)} \comp R = R = R \comp \diag{\Cod(R)}$.
	\end{frame}
	
	\begin{frame}{Diagrama de uma relação finita}
		\passo
		Uma relação $R$ pode ser representada por gráfico cartesiano, com $\Dom(R)$ nas colunas e
		$\Cod(R)$ nas linhas.
		Para $\Dom(R) = \{1,2,3\}$, $\Cod(R) = \{a,b,c\}$ e $R = \{(1,b),(2,a),(3,b)\}$:
		\begin{center}
			\uncover<+->{$R = {}$\relmat[\relaxes]{$c$,$b$,$a$}{1,2,3}{{0,0,0},{1,0,1},{0,1,0}}}
			\qquad
			\uncover<+->{$\inv{R} = {}$\relmat{3,2,1}{$a$,$b$,$c$}{{0,1,0},{1,0,0},{0,1,0}}}
		\end{center}
		
		\medskip
		\passo
		Cada definição tem uma leitura no diagrama:
		\begin{itemize}
			\item<+-> $\inv{R}$ permuta os eixos,
			\item<+-> $\Img(R)$ é o conjunto dos rótulos das linhas não vazias,
			\item<+-> $\Fix(R)$ é o conjunto dos elementos sobre a diagonal,
			\item<+-> o desenho de $R \comp S$ é o produto booleano do desenho de $S$ pelo de $R$.
		\end{itemize}
	\end{frame}
	
	\begin{frame}{Restrição funcional}
		\passo
		De agora em diante $R$ é uma função. Da definição extraímos os dois únicos instrumentos
		usados adiante:
		\begin{itemize}
			\item<+-> \textbf{existência}: para todo $a \in \Dom(R)$ há $b \in \Cod(R)$ com $aRb$,
			\item<+-> \textbf{unicidade}: se $aRb$ e $aRc$, então $b = c$.
		\end{itemize}
		
		\medskip
		\passo
		Juntos, eles dizem que cada coluna do desenho tem exatamente um elemento, que é a leitura
		usual de gráfico. Tome $\Dom(R) = \{1,2\}$ e $\Cod(R) = \{a,b,c\}$:
		\begin{center}
			\begin{tabular}{ccc}
				\uncover<+->{\begin{tabular}{@{}c@{}}
						\relmat[\hlcol{1}]{$c$,$b$,$a$}{1,2}{{0,0},{0,0},{0,1}} \\[4pt]
						\small falha a existência
				\end{tabular}} &
				\uncover<+->{\begin{tabular}{@{}c@{}}
						\relmat[\hlcol{1}]{$c$,$b$,$a$}{1,2}{{0,1},{1,0},{1,0}} \\[4pt]
						\small falha a unicidade
				\end{tabular}} &
				\uncover<+->{\begin{tabular}{@{}c@{}}
						\relmat{$c$,$b$,$a$}{1,2}{{0,0},{1,0},{0,1}} \\[4pt]
						\small função
				\end{tabular}}
			\end{tabular}
		\end{center}
		
		\medskip
		\passo
		Toda demonstração adiante consiste em encontrar o ponto certo para aplicar a unicidade.
	\end{frame}
	
	\begin{frame}{Definições algébricas das propriedades}
		\passo
		\begin{center}
			\begin{tabular}{ll}
				\multicolumn{2}{l}{\textbf{Tipo heterogêneo} (fazem sentido com $\Dom(R) \neq \Cod(R)$)}\\[2pt]
				total (à esquerda) & $\diag{\Dom(R)} \subseteq R \comp \inv{R}$ \\
				injetiva       & $R \comp \inv{R} \subseteq \diag{\Dom(R)}$ \\
				sobrejetiva    & $\diag{\Cod(R)} \subseteq \inv{R} \comp R$ \\
				unívoca        & $\inv{R} \comp R \subseteq \diag{\Cod(R)}$ \\[8pt]
				\multicolumn{2}{l}{\textbf{Tipo homogêneo}}\\[2pt]
				reflexiva      & $\diag{\Dom(R)} \subseteq R$ \\
				simétrica      & $\inv{R} \subseteq R$ \\
				anti-simétrica & $R \cap \inv{R} \subseteq \diag{\Dom(R)}$ \\
				transitiva     & $R \comp R \subseteq R$
			\end{tabular}
		\end{center}
		
		\medskip
		\passo
		Uma relação é função quando é total e unívoca.
	\end{frame}
	
	\begin{frame}{O que será provado}
		\passo
		Para $R$ função, as quatro propriedades homogêneas se reduzem a
		\begin{center}
			\begin{tabular}{ll}
				reflexiva      & $\Dom(R) \subseteq \Fix(R)$ \\
				simétrica      & $\Dom(R) \subseteq \Fix(R \comp R)$ \\
				anti-simétrica & $\Fix(R \comp R) \subseteq \Fix(R)$ \\
				transitiva     & $\Img(R) \cap \Dom(R) \subseteq \Fix(R)$
			\end{tabular}
		\end{center}
		
		\medskip
		\passo
		As quatro comparam termos de uma mesma cadeia, que vale para qualquer relação:
		\[
		\Fix(R) \subseteq \Fix(R \comp R) \subseteq \Img(R) \cap \Dom(R) \subseteq \Dom(R).
		\]
		\passo
		A primeira inclusão é $aRa \Rightarrow a\,(R \comp R)\,a$. Na segunda, de $aRb$ e $bRa$ vem
		$a \in \Dom(R)$ e $a \in \Img(R)$. A terceira é imediata.
		
		\medskip
		\passo
		Cada propriedade colapsa um trecho da cadeia. Reflexividade colapsa tudo e dá
		$R = \diag{\Dom(R)}$, simetria colapsa os três termos externos e torna $R$ uma involução de
		$\Dom(R)$. As duas últimas são bem mais fracas, e só dizem algo sobre $\Img(R) \cap \Dom(R)$.
	\end{frame}
	
	\begin{frame}{Composição com a inversa}
		\passo
		\begin{align*}
			a\,(R \comp \inv{R})\,a' &\iff \exists b,\ (aRb \wedge a'Rb), \\
			b\,(\inv{R} \comp R)\,b' &\iff \exists a,\ (aRb \wedge aRb').
		\end{align*}
		\passo
		A primeira relaciona duas colunas que dividem uma linha, a segunda duas linhas que dividem
		uma coluna.
		
		\medskip
		\passo
		Totalidade e injetividade são condições sobre $R \comp \inv{R}$. A primeira pede ao menos
		um elemento em cada coluna, a segunda pede que duas colunas nunca dividam uma linha, isto é,
		no máximo um elemento por linha. Sobrejetividade e univocidade são as condições espelhadas
		sobre $\inv{R} \comp R$.
		
		\medskip
		\passo
		Para $R$ função, $R \comp \inv{R}$ é a relação \emph{ter a mesma imagem}, e
		$\inv{R} \comp R = \diag{\Img(R)}$:
		\begin{center}
			\uncover<+->{$R$\,\relmat{$c$,$b$,$a$}{1,2,3}{{0,0,0},{1,0,1},{0,1,0}}}
			\;
			\uncover<+->{$R \comp \inv{R}$\,\relmat[\reldiag]{3,2,1}{1,2,3}{{1,0,1},{0,1,0},{1,0,1}}}
			\;
			\uncover<+->{$\inv{R} \comp R$\,\relmat[\reldiag]{$c$,$b$,$a$}{$a$,$b$,$c$}{{0,0,0},{0,1,0},{1,0,0}}}
		\end{center}
		\onslide<.->
		\ordemnota
	\end{frame}
	
	\begin{frame}{Composição com ela mesma}
		\passo
		\[
		a\,(R \comp R)\,c \iff \exists b,\ (aRb \wedge bRc).
		\]
		\passo
		São os pares ligados por um caminho de dois passos. O ponto intermediário precisa ser ao
		mesmo tempo imagem e ponto de partida, logo $R \comp R$ só tem conteúdo sobre
		$\Img(R) \cap \Dom(R)$, e é vazia quando $\Dom(R)$ e $\Cod(R)$ são disjuntos.
		
		\medskip
		\passo
		Para $R$ função, $R \comp R$ é aplicar $R$ duas vezes, e é unívoca: a unicidade em $a$ fixa
		o intermediário, e a unicidade nele fixa o resultado. Todos os seus pares partem de pontos
		de $\Dom(R)$.
		
		\medskip
		\passo
		Com $\Dom(R) = \Cod(R) = \{1,2,3\}$ e $R = \{(1,2),(2,3),(3,3)\}$, dois passos levam tudo
		para o ponto fixo $3$. Como $(1,3) \notin R$, essa função não é transitiva:
		\begin{center}
			\uncover<+->{$R$\,\relmat[\reldiag]{3,2,1}{1,2,3}{{0,1,1},{1,0,0},{0,0,0}}}
			\qquad
			\uncover<+->{$R \comp R$\,\relmat[\reldiag]{3,2,1}{1,2,3}{{1,1,1},{0,0,0},{0,0,0}}}
		\end{center}
		\onslide<.->
		\ordemnota
	\end{frame}
	
	\begin{frame}{Reflexividade}
		\passo
		\begin{prop}\label{prop:refl}
			São equivalentes: $R$ é reflexiva, $\Dom(R) \subseteq \Fix(R)$, $R = \diag{\Dom(R)}$.
			Nesse caso $\Dom(R) \subseteq \Cod(R)$.
		\end{prop}
		
		\passo
		\begin{proof}
			As duas primeiras são a mesma frase, porque $\diag{\Dom(R)} \subseteq R$ diz que $aRa$
			vale para todo $a \in \Dom(R)$. Delas para a terceira: se $aRb$, a unicidade em $a$,
			comparada com $aRa$, dá $b = a$, donde $R \subseteq \diag{\Dom(R)}$, e a inclusão
			contrária é a própria reflexividade. A volta é imediata. Por fim, de $aRa$ vem
			$a \in \Cod(R)$.
		\end{proof}
		
		\medskip
		\begin{columns}[c]
			\begin{column}{0.68\textwidth}
				\uncover<+->{Uma função reflexiva é a inclusão $\Dom(R) \hookrightarrow \Cod(R)$. Em
					particular ela é injetiva, e sobrejetiva exatamente quando $\Dom(R) = \Cod(R)$.}
			\end{column}
			\begin{column}{0.30\textwidth}
				\begin{center}
					\uncover<+->{\begin{tabular}{@{}c@{}}
							\relmat[\hlrow{1}]{3,2,1}{1,2}{{0,0},{0,1},{1,0}} \\[4pt]
							\scriptsize $\Dom(R) = \{1,2\}$ \\
							\scriptsize $\Cod(R) = \{1,2,3\}$
					\end{tabular}}
				\end{center}
			\end{column}
		\end{columns}
	\end{frame}
	
	\mode<handout>{
		\begin{frame}{Reflexividade, demonstração completa}
			A equivalência entre $\diag{\Dom(R)} \subseteq R$ e $\Dom(R) \subseteq \Fix(R)$ é a
			definição de $\Fix$, porque ambas dizem que $aRa$ vale para todo $a \in \Dom(R)$.
			
			\medskip
			Suponha $R$ reflexiva e seja $aRb$. Como também $aRa$, a unicidade em $a$ dá $b = a$, logo
			$(a,b) \in \diag{\Dom(R)}$ e $R \subseteq \diag{\Dom(R)}$. A inclusão
			$\diag{\Dom(R)} \subseteq R$ é a hipótese, portanto $R = \diag{\Dom(R)}$.
			
			\medskip
			Reciprocamente, se $R = \diag{\Dom(R)}$, então $\diag{\Dom(R)} \subseteq R$ e $R$ é
			reflexiva.
			
			\medskip
			Para a última afirmação, dado $a \in \Dom(R)$ vale $aRa$, e como
			$R \subseteq \Dom(R) \times \Cod(R)$ isso dá $a \in \Cod(R)$. Logo
			$\Dom(R) \subseteq \Cod(R)$.
		\end{frame}
	}
	
	\begin{frame}{Simetria}
		\passo
		\begin{prop}\label{prop:sim}
			São equivalentes: (i) $R$ é simétrica, (ii) $\Dom(R) \subseteq \Fix(R \comp R)$,
			(iii) $R \comp R = \diag{\Dom(R)}$.
			Nesse caso $\inv{R} = R$, $\Img(R) \subseteq \Dom(R) \subseteq \Cod(R)$, e $R$ é uma
			involução de $\Dom(R)$.
		\end{prop}
		
		\passo
		\begin{proof}
			(i)\,$\Rightarrow$\,(ii): dado $a \in \Dom(R)$, tome $b$ com $aRb$, e de $bRa$ vem
			$a\,(R \comp R)\,a$.
			(ii)\,$\Rightarrow$\,(iii): $R \comp R$ é unívoca e parte de $\Dom(R)$, logo um ponto
			fixo em cada $a$ esgota a relação.
			(iii)\,$\Rightarrow$\,(i): de $aRb$ e $a\,(R \comp R)\,a$, a unicidade em $a$ identifica
			o intermediário com $b$.
		\end{proof}
		
		\begin{columns}[c]
			\begin{column}{0.46\textwidth}
				\uncover<+->{\footnotesize O desenho é uma matriz de permutação involutiva, simétrica pela diagonal.
					A transposição de $1$ e $2$ mostra que simetria não força $R = \diag{\Dom(R)}$.\par}
			\end{column}
			\begin{column}{0.52\textwidth}
				\begin{center}
					\uncover<+->{$R$\,\relmat[\reldiag]{3,2,1}{1,2,3}{{0,0,1},{1,0,0},{0,1,0}}}
					\quad
					\uncover<+->{$R \comp R$\,\relmat[\reldiag]{3,2,1}{1,2,3}{{0,0,1},{0,1,0},{1,0,0}}}
				\end{center}
			\end{column}
		\end{columns}
	\end{frame}
	
	\mode<handout>{
		\begin{frame}{Simetria, demonstração completa}
			Suponha $R$ simétrica. Dado $a \in \Dom(R)$, tome $b$ com $aRb$ pela existência. A
			simetria dá $bRa$, logo $a\,(R \comp R)\,a$ e $\Dom(R) \subseteq \Fix(R \comp R)$.
			
			\medskip
			Suponha $\Dom(R) \subseteq \Fix(R \comp R)$ e seja $a\,(R \comp R)\,c$, digamos $aRb$ e
			$bRc$. Então $a \in \Dom(R)$, logo há $b'$ com $aRb'$ e $b'Ra$. A unicidade em $a$ dá
			$b' = b$ e a unicidade em $b$ dá $c = a$. Junto com a hipótese, isso é
			$R \comp R = \diag{\Dom(R)}$.
			
			\medskip
			Suponha $R \comp R = \diag{\Dom(R)}$ e seja $aRb$. De $a\,(R \comp R)\,a$ há $b'$ com
			$aRb'$ e $b'Ra$, e a unicidade em $a$ dá $b' = b$, donde $bRa$ e $R$ é simétrica.
			
			\medskip
			Para o resto, de $aRb$ vem $bRa$, logo $b \in \Dom(R)$ e $a \in \Img(R) \subseteq \Cod(R)$,
			o que dá $\Img(R) \subseteq \Dom(R) \subseteq \Cod(R)$. Invertendo $\inv{R} \subseteq R$
			obtemos $R \subseteq \inv{R}$, logo $\inv{R} = R$, e $R \comp R = \diag{\Dom(R)}$ faz de
			$R$ uma involução de $\Dom(R)$, em particular uma bijeção de $\Dom(R)$ sobre $\Dom(R)$.
		\end{frame}
	}
	
	\begin{frame}{Anti-simetria}
		\passo
		É a única das quatro que usa $R \cap \inv{R}$, os pares recíprocos:
		\[
		a\,(R \cap \inv{R})\,b \iff (aRb \wedge bRa).
		\]
		\passo
		No desenho, é o que sobrevive à interseção de $R$ com sua reflexão pela diagonal. Ela sempre
		contém $\diag{\Fix(R)}$, e a condição pede que nada mais sobreviva.
		
		\passo
		\begin{prop}\label{prop:anti}
			$R$ é anti-simétrica se, e só se, $\Fix(R \comp R) \subseteq \Fix(R)$.
		\end{prop}
		
		\passo
		\begin{proof}
			\volta Sejam $aRb$ e $bRa$. Então $a \in \Fix(R \comp R) \subseteq \Fix(R)$, logo $aRa$, e
			a unicidade em $a$ dá $b = a$.
			\ida Seja $a \in \Fix(R \comp R)$, digamos $aRb$ e $bRa$. A anti-simetria dá $a = b$,
			donde $aRa$.
		\end{proof}
		
		\begin{columns}[c]
			\begin{column}{0.62\textwidth}
				\uncover<+->{\small A condição proíbe $a \neq b$ com $aRb$ e $bRa$. Ela vale vacuamente quando
					$\Img(R) \cap \Dom(R) = \varnothing$.\par}
			\end{column}
			\begin{column}{0.36\textwidth}
				\begin{center}
					\begin{tabular}{cc}
						\uncover<+->{\begin{tabular}{@{}c@{}}
								\relmat[\reldiag\hlcell{1}{1}\hlcell{2}{2}]{2,1}{1,2}{{1,0},{0,1}} \\[3pt]
								\scriptsize proibido
						\end{tabular}} &
						\uncover<+->{\begin{tabular}{@{}c@{}}
								\relmat[\reldiag]{2,1}{1,2}{{0,0},{1,1}} \\[3pt]
								\scriptsize permitido
						\end{tabular}}
					\end{tabular}
				\end{center}
			\end{column}
		\end{columns}
	\end{frame}
	
	\mode<handout>{
		\begin{frame}{Anti-simetria, demonstração completa}
			Suponha $\Fix(R \comp R) \subseteq \Fix(R)$ e sejam $aRb$ e $bRa$. Então
			$a\,(R \comp R)\,a$, logo $a \in \Fix(R \comp R)$ e portanto $aRa$. A unicidade em $a$,
			comparando $aRb$ com $aRa$, dá $b = a$. Assim $R \cap \inv{R} \subseteq \diag{\Dom(R)}$.
			
			\medskip
			Reciprocamente, suponha $R$ anti-simétrica e seja $a \in \Fix(R \comp R)$. Há $b$ com
			$aRb$ e $bRa$, logo $(a,b) \in R \cap \inv{R} \subseteq \diag{\Dom(R)}$, donde $a = b$ e
			$aRa$.
			
			\medskip
			Vale sempre $\diag{\Fix(R)} \subseteq R \cap \inv{R}$, porque $aRa$ dá o par recíproco
			trivial. Sob anti-simetria a inclusão contrária também vale, e as duas juntas dão
			$R \cap \inv{R} = \diag{\Fix(R)}$.
		\end{frame}
	}
	
	\begin{frame}{Transitividade}
		\passo
		\begin{prop}\label{prop:trans}
			$R$ é transitiva se, e só se, $\Img(R) \cap \Dom(R) \subseteq \Fix(R)$.
			Se além disso $\Img(R) \subseteq \Dom(R)$, isto equivale a $R \comp R = R$.
		\end{prop}
		
		\passo
		\begin{proof}
			\ida De $aRb$ e $bRc$, a transitividade dá $aRc$, e a unicidade em $a$ dá $c = b$, donde
			$bRb$.
			\volta De $aRb$ e $bRc$ vem $b \in \Img(R) \cap \Dom(R)$, logo $bRb$, e a unicidade em $b$
			dá $c = b$, donde $aRc$.
		\end{proof}
		
		\medskip
		\passo
		A cadeia dá $\Fix(R) \subseteq \Img(R) \cap \Dom(R)$ sem hipótese alguma, logo sob
		transitividade os elementos de $\Img(R) \cap \Dom(R)$ são exatamente os pontos fixos. No caso
		$\Img(R) \subseteq \Dom(R)$, as funções transitivas são as idempotentes, isto é, as retrações
		de $\Dom(R)$ sobre $\Fix(R)$.
	\end{frame}
	
	\mode<handout>{
		\begin{frame}{Transitividade, demonstração completa}
			Seja $R$ transitiva e $b \in \Img(R) \cap \Dom(R)$. Tome $a$ com $aRb$ e, pela existência,
			$c$ com $bRc$. A transitividade dá $aRc$ e a unicidade em $a$ dá $c = b$, logo $bRb$.
			
			\medskip
			Reciprocamente, suponha $\Img(R) \cap \Dom(R) \subseteq \Fix(R)$ e sejam $aRb$ e $bRc$.
			Então $b \in \Img(R) \cap \Dom(R)$, logo $bRb$, e a unicidade em $b$ dá $c = b$, donde
			$aRc$.
			
			\medskip
			Para a última afirmação, $R \comp R \subseteq R$ é a transitividade. Se $aRb$, então
			$b \in \Dom(R)$ por hipótese, logo $bRb$ pelo que acabamos de ver, e portanto
			$(a,b) \in R \comp R$. Assim $R \subseteq R \comp R$ e vale a igualdade.
		\end{frame}
	}
	
	\begin{frame}{Diagrama da transitividade}
		\passo
		A relação $R = \{(1,2),(2,3),(3,1)\}$ falha, porque os desenhos de $R$ e de
		$R \comp R$ são disjuntos. Nenhum elemento cai sobre a diagonal.
		\begin{center}
			\uncover<+->{$R$\,\relmat[\reldiag]{3,2,1}{1,2,3}{{0,1,0},{1,0,0},{0,0,1}}}
			\qquad
			\uncover<+->{$R \comp R$\,\relmat[\reldiag]{3,2,1}{1,2,3}{{1,0,0},{0,0,1},{0,1,0}}}
		\end{center}
		
		\medskip
		\passo
		Já $R = \{(1,1),(2,2),(3,1)\}$ é idempotente, com $\Fix(R) = \{1,2\}$. Os elementos ficam
		confinados às linhas de $\Fix(R)$, e a coluna de cada ponto fora de $\Fix(R)$ aponta
		para uma dessas linhas.
		\begin{center}
			\uncover<+->{$R$\,\relmat[\reldiag]{3,2,1}{1,2,3}{{0,0,0},{0,1,0},{1,0,1}}}
			\qquad
			\uncover<+->{$R \comp R$\,\relmat[\reldiag]{3,2,1}{1,2,3}{{0,0,0},{0,1,0},{1,0,1}}}
		\end{center}
	\end{frame}
	
	\begin{frame}{Resumo}
		\passo
		Cada propriedade colapsa um trecho da cadeia.
		
		\begin{columns}[c]
			\begin{column}{0.42\textwidth}
				\begin{center}\uncover<+->{\scalebox{0.82}{\cadeia}}\end{center}
			\end{column}
			\begin{column}{0.56\textwidth}
				\small
				\uncover<+->{\begin{tabular}{ll}
						reflexiva & colapsa os quatro \\
						simétrica & colapsa os três externos \\
						anti-simétrica & colapsa os dois internos \\
						transitiva & colapsa os três internos
				\end{tabular}}
				
				\smallskip
				\uncover<+->{As implicações se leem aqui. Reflexiva implica as outras três, e transitiva implica
					anti-simétrica. Simétrica com anti-simétrica, ou com transitiva, devolve a
					reflexividade.\par}
			\end{column}
		\end{columns}
		
		\smallskip
		\passo
		Os exemplos abaixo refutam as demais implicações:
		\begin{center}
			\uncover<+->{\scriptsize\setlength{\tabcolsep}{2pt}%
				\begin{tabular}{lcccc}
					& \relgrid[\reldiag]{{0,1},{1,0}} & \relgrid[\reldiag]{{1,0},{0,1}} &
					\relgrid[\reldiag]{{0,0},{1,1}} & \relgrid[\reldiag]{{0,1,0},{1,0,0},{0,0,1}} \\[2pt]
					& $\{(1,1),(2,2)\}$ & $\{(1,2),(2,1)\}$ & $\{(1,1),(2,1)\}$ & $\{(1,2),(2,3),(3,1)\}$ \\ \hline
					reflexiva & \checkmark & & & \\
					simétrica & \checkmark & \checkmark & & \\
					anti-simétrica & \checkmark & & \checkmark & \checkmark \\
					transitiva & \checkmark & & \checkmark & \\
			\end{tabular}}
		\end{center}
		\onslide<.->
		{\footnotesize As três primeiras em $\Dom(R) = \Cod(R) = \{1,2\}$, a última em
			$\Dom(R) = \Cod(R) = \{1,2,3\}$.}
	\end{frame}
	
	\begin{frame}{Ordem e equivalência}
		\passo
		Aqui, \emph{ordem parcial} significa reflexiva, anti-simétrica e transitiva, e
		\emph{relação de equivalência} significa reflexiva, simétrica e transitiva. Nenhuma das
		duas supõe $\Dom(R) = \Cod(R)$.
		
		\passo
		\begin{cor}
			São equivalentes: $R$ é ordem parcial, $R$ é equivalência, $R = \diag{\Dom(R)}$.
		\end{cor}
		
		\passo
		\begin{proof}
			Ambas incluem reflexividade, que dá $R = \diag{\Dom(R)}$ pela
			Proposição~\ref{prop:refl}, e $\diag{\Dom(R)}$ é ordem parcial e equivalência.
		\end{proof}
		
		\medskip
		\begin{columns}[c]
			\begin{column}{0.58\textwidth}
				\uncover<+->{A igualdade é a única relação de ordem que também é função, correspondendo à ordem discreta. É também a única relação de equivalência que é função, correspondendo à partição em singletons. Ordens e equivalências não triviais não são unívocas, e essa falha aparece nas colunas destacadas.\par}
			\end{column}
			\begin{column}{0.40\textwidth}
				\begin{center}
					\begin{tabular}{cc}
						\uncover<+->{\begin{tabular}{@{}c@{}}
								\relmat[\reldiag\hlcol{1}]{3,2,1}{1,2,3}{{1,1,1},{1,1,0},{1,0,0}} \\[3pt]
								\scriptsize ordem $\leq$
						\end{tabular}} &
						\uncover<+->{\begin{tabular}{@{}c@{}}
								\relmat[\reldiag\hlcol{1}]{3,2,1}{1,2,3}{{0,0,1},{1,1,0},{1,1,0}} \\[3pt]
								\scriptsize equivalência
						\end{tabular}}
					\end{tabular}
				\end{center}
			\end{column}
		\end{columns}
	\end{frame}
	
	\begin{frame}{Injetividade e sobrejetividade}
		\passo
		Reflexividade e simetria forçam $\Img(R) \subseteq \Dom(R) \subseteq \Cod(R)$ e implicam
		injetividade. Anti-simetria e transitividade valem vacuamente quando
		$\Img(R) \cap \Dom(R) = \varnothing$, e não implicam nenhuma das duas.
		
		\passo
		\begin{cor}
			\begin{enumerate}
				\item<+-> Reflexiva $\Rightarrow$ injetiva, e sobrejetiva se, e só se, $\Dom(R) = \Cod(R)$.
				\item<+-> Simétrica $\Rightarrow$ injetiva, com $R$ bijeção de $\Dom(R)$ sobre $\Dom(R)$.
				\item<+-> Simétrica e anti-simétrica, ou simétrica e transitiva, $\Rightarrow R = \diag{\Dom(R)}$.
				\item<+-> Transitiva e injetiva, com $\Img(R) \subseteq \Dom(R)$, $\Rightarrow R = \diag{\Dom(R)}$.
			\end{enumerate}
		\end{cor}
		
		\passo
		\begin{proof}
			(1) é a Proposição~\ref{prop:refl}. Em (2), $\inv{R} = R$ dá
			$R \comp \inv{R} = R \comp R = \diag{\Dom(R)}$.
			Em (3), a simetria dá $\Dom(R) \subseteq \Fix(R \comp R)$, e tanto a anti-simetria quanto
			a transitividade, esta via $\Fix(R \comp R) \subseteq \Img(R) \cap \Dom(R)$, dão
			$\Fix(R \comp R) \subseteq \Fix(R)$, donde $\Dom(R) \subseteq \Fix(R)$.
			Em (4), de $aRb$ com $b \in \Dom(R)$ vem $bRb$, e a injetividade dá $a = b$.
		\end{proof}
	\end{frame}
	
\end{document}